%  =============================================================================
%  Datei-Hinweise & Konfiguration: siehe config.md
%  -> Abschnitt "anhang/quellcodearchiv.tex"
%
%  PFLEGE / REGENERIERUNG (z. B. nach der Echtgeld-Woche):
%    1. Im Projektordner archiv_manifest.csv anpassen (neue Datei = neue
%       Zeile; Status "geplant" -> "aktiv" setzen).
%    2. python3 build_quellcode_archiv.py --version V3 --datum <TTMONJJJJ>
%    3. Die vier GENERIERTEN Dateien archiv_meta.tex, archiv_quellcode.tex,
%       archiv_tests.tex, archiv_daten.tex aus archiv_tex/ nach Overleaf in
%       den Ordner anhang/ hochladen (ersetzen). Diese Datei hier bleibt
%       unveraendert.
%  =============================================================================

\chapter{Anhang: Projektstruktur und Quellcodearchiv}
\label{cha:anhang-archiv}

% Generierte Metadaten des Archivs (Name, Datum, SHA-256, Umfang).
%% ---------------------------------------------------------------------------
%% GENERIERT durch build_quellcode_archiv.py am 27.07.2026 12:00 — NICHT von Hand
%% pflegen; Aenderungen in archiv_manifest.csv vornehmen und Skript ausfuehren.
%% ---------------------------------------------------------------------------
\newcommand{\QuellcodeArchivName}{FOM\_Projektarbeit\_KI\_Quellcode\_FischerKiss\_27JUL2026\_V3.zip}
\newcommand{\QuellcodeArchivDatum}{27.07.2026}
\newcommand{\QuellcodeArchivVersion}{V3}
\newcommand{\QuellcodeArchivSHA}{2c56790b80c2f793fa47665b041840a27aefa26596cb04ea80d4560772d5a09e}
\newcommand{\QuellcodeArchivDateien}{112}
\newcommand{\QuellcodeArchivGroesse}{0.7}


Dieser Anhang dokumentiert den Quellcodestand des in \cref{cha:umsetzung}
beschriebenen Handelsagenten. Der vollständige Quellcode wird nicht im Anhang
abgedruckt, sondern liegt der Arbeit als elektronisches Archiv bei; die
nachfolgenden Abschnitte weisen dessen Identität, Aufbau und Inhalt aus. Damit
bleibt jede in \cref{sec:umsetzung-implementierung} beschriebene Absicherung im
Quelltext nachvollziehbar, ohne den Anhang mit rund 7.000 Zeilen Quelltext zu
belasten.

\section{Projektstruktur}
\label{sec:anhang-projektstruktur}

\Cref{fig:projektstruktur} zeigt die Struktur des Projekts. Es handelt sich um
einen einzigen Codestand, der auf beiden virtuellen Maschinen ausgerollt wird;
die Zuordnung der Module zu collector-vm und agent-vm ist je Datei vermerkt.
Die Zuordnung jeder Datei zu ihrem Zweck und ihrem Ort im Archiv dokumentieren
die Tabellen der \cref{sec:anhang-archiv-quellcode,sec:anhang-archiv-tests,sec:anhang-archiv-daten}.

\begin{figure}[htbp]
    \centering
    \includesvg[width=0.92\textwidth,height=0.88\textheight,keepaspectratio]{projektstruktur}
    \caption[Projektstruktur des Handelsagenten]{Projektstruktur des
        Handelsagenten: ein Codestand für beide virtuellen Maschinen
        (C = collector-vm, A = agent-vm).\quelleeigenki{Anthropic Claude Fable~5 (Anthropic, 2026)}}
    \label{fig:projektstruktur}
\end{figure}

\section{Archividentität und Redaktion}
\label{sec:anhang-archiv-identitaet}

Das Archiv ist über Dateinamen, Erstellungsdatum und Prüfsumme eindeutig
identifiziert; die Prüfsumme erlaubt den Nachweis, dass der beigelegte Stand
dem hier dokumentierten entspricht.

\begin{table}[htbp]
  \centering
  \begin{tabular}{l l}
    \hline
    Merkmal & Wert \\
    \hline
    Dateiname & \texttt{\footnotesize\QuellcodeArchivName} \\
    Version / Datum & \QuellcodeArchivVersion{} / \QuellcodeArchivDatum \\
    Umfang & \QuellcodeArchivDateien{} Dateien, \QuellcodeArchivGroesse{}\,MB \\
    Integrität (SHA-256) & \texttt{\footnotesize\QuellcodeArchivSHA} \\
    \hline
  \end{tabular}
  \caption[Identität des Quellcodearchivs]{Identität des beigelegten
    Quellcodearchivs.\quelleeigen}
  \label{tab:anhang-archiv-identitaet}
\end{table}

Für die Beilage gelten zwei redaktionelle Regeln, die die Programmlogik
unberührt lassen: Die Rufnummern des Messaging-Kanals sind redigiert
(\texttt{+49XXXXXXXXXXX}), und Zugangsdaten erscheinen an keiner Stelle, da sie
ausschließlich im verschlüsselten Geheimnisspeicher liegen -- die Datei
\texttt{agent.env} und der Inhalt von \texttt{\textasciitilde/.fom-agent/}
sind deshalb nicht Bestandteil des Archivs. Die unverändert eingebundene
Fremdbibliothek Chart.js (\texttt{webfeed/app/static/chart.umd.js}, MIT-Lizenz)
ist als Drittcode ebenfalls nicht beigelegt; im Übrigen liegt der Quellcode
unverändert und lauffähig vor.

\section{Aufbau des Archivs}
\label{sec:anhang-archiv-aufbau}

Das Archiv gliedert sich in drei Basisordner sowie eine erläuternde
\texttt{README.md}: Der Ordner \texttt{Quellcode/} enthält sämtliche Module,
Skripte und Konfigurationsdateien in der flachen Struktur der
Produktivumgebung (vgl. \cref{fig:projektstruktur}), wodurch der Code ohne
Umbau lauffähig bleibt; der Ordner \texttt{Tests/} enthält die
Offline-Testsuiten einschließlich des Testtreibers \texttt{run\_tests.sh}; der
Ordner \texttt{Daten/} enthält die Daten- und Evaluationsstände mit
ausgewiesenem Snapshot-Datum.

\section{Quellcode}
\label{sec:anhang-archiv-quellcode}

%% ---------------------------------------------------------------------------
%% GENERIERT durch build_quellcode_archiv.py am 25.07.2026 08:59 — NICHT von Hand
%% pflegen; Aenderungen in archiv_manifest.csv vornehmen und Skript ausfuehren.
%% ---------------------------------------------------------------------------
\begin{longtable}{@{}p{0.29\textwidth}p{0.155\textwidth}p{0.455\textwidth}@{}}
  \caption[Quellcodedateien des Archivs]{Ordner \texttt{Quellcode/} des Archivs: Dateien und Zwecke, gegliedert nach Funktionsbereichen der Projektstruktur.\quelleeigen}\label{tab:anhang-archiv-quellcode}\\
  \toprule
  Datei & Typ & Zweck \\
  \midrule
  \endfirsthead
  \multicolumn{3}{@{}l}{\itshape Fortsetzung von vorheriger Seite}\\
  \toprule
  Datei & Typ & Zweck \\
  \midrule
  \endhead
  \midrule
  \multicolumn{3}{r@{}}{\itshape Fortsetzung auf der n\"achsten Seite}\\
  \endfoot
  \bottomrule
  \endlastfoot
  \multicolumn{3}{@{}l}{\textbf{Konfiguration}}\\*[2pt]
  \texttt{agentconfig.yaml} & Kon\-fi\-gu\-ra\-tion & Betriebskonfiguration: Strategien, Deployments, Sprachmodelle, Zeitbudgets, ML-Prüfinstanz \\
  \texttt{wknassign.yaml} & Kon\-fi\-gu\-ra\-tion & Zentrale Produktzuordnung Strategie zu Handelsinstrument (WKN/ISIN) \\
  \addlinespace[6pt]
  \multicolumn{3}{@{}l}{\textbf{Datenerfassung (collector-vm)}}\\*[2pt]
  \texttt{datacollect.py} & Python & Kontinuierliche Marktdatenerfassung mit Sanitätsfilter, Bar-Verdichtung und FRED-Zweitquellenprüfung \\
  \texttt{market\_\allowbreak calendar.py} & Python & Börsenkalender: Handelstage, Feiertage, Sitzungsfenster \\
  \texttt{collector\_\allowbreak summary.py} & Python & Tageszusammenfassung der Datenerfassung für den Berichtsversand \\
  \addlinespace[6pt]
  \multicolumn{3}{@{}l}{\textbf{Strategieschicht}}\\*[2pt]
  \texttt{strategy.py} & Python & Strategie-Handler: Signalberechnung, DecisionRequest, Audit-Flags, Kostenkontext, Instrumentenkatalog, Backtest-Kommando \\
  \texttt{strategyloader.py} & Python & Plugin-Lader: aktive Strategien, Deployments, ML-Konfiguration \\
  \texttt{plstrategy\_\allowbreak sma.py} & Python & Plugin der SMA-Regimestrategie (Kreuzungssignal) \\
  \texttt{plstrategy\_\allowbreak smacross.py} & Python & Plugin-Variante des Kreuzungssignals \\
  \texttt{plstrategy\_\allowbreak smatrend.py} & Python & Trendfolge-Alternative für den A/B-Vergleich \\
  \texttt{plstrategy\_\allowbreak buyhold.py} & Python & Vergleichsmaßstab passives Halten \\
  \texttt{plstrategy\_\allowbreak momentum.py} & Python & Vergleichsmaßstab Momentum \\
  \texttt{wknassign\_\allowbreak sma.py} & Python & Veralteter Platzhalter, ersetzt durch wknassign.yaml (dokumentiert die Migration) \\
  \texttt{wknassign\_\allowbreak smacross.py} & Python & Veralteter Platzhalter, ersetzt durch wknassign.yaml \\
  \texttt{wknassign\_\allowbreak momentum.py} & Python & Veralteter Platzhalter, ersetzt durch wknassign.yaml \\
  \addlinespace[6pt]
  \multicolumn{3}{@{}l}{\textbf{Prüfinstanzen}}\\*[2pt]
  \texttt{auditor.py} & Python & Multi-LLM-Voting: Pre-Gate, unabhängige Voten, strikte Mehrheit, fail-safe Validator \\
  \texttt{mlforecast.py} & Python & Statistische Prüfinstanz: logistische Regression, Walk-Forward-Prüfung, Schatten- und Wirkmodus \\
  \addlinespace[6pt]
  \multicolumn{3}{@{}l}{\textbf{Orderausführung und Bankanbindung}}\\*[2pt]
  \texttt{orchestrate.py} & Python & Abgesicherte Orderausführung: Frische-Gates, Intent-Journal, Orderbuch-Reconciliation, Circuit-Breaker \\
  \texttt{broker.py} & Python & comdirect-Adapter: OAuth2, Session-TAN, Depot- und Orderbuchzugriff, tagesgültige Limit-Orders \\
  \texttt{credstore.py} & Python & Verschlüsselte Ablage der Zugangsdaten, getrennt je Konto \\
  \texttt{config.py} & Python & Lader der Umgebungskonfiguration \\
  \addlinespace[6pt]
  \multicolumn{3}{@{}l}{\textbf{Betrieb und Steuerung}}\\*[2pt]
  \texttt{run\_\allowbreak cycle.sh} & Shell & Täglicher Entscheidungszyklus: Preflight, Strategie, Audit, Ausführung, Bericht \\
  \texttt{run\_\allowbreak all.sh} & Shell & Cron-Einstieg: Entscheidungszyklus je Konto/Deployment \\
  \texttt{control.py} & Python & Befehlsprotokoll des Steuerkanals: Not-Aus, Broker-Login, ML-Umschaltung, Status \\
  \texttt{signal\_\allowbreak dispatcher.py} & Python & Empfangsschleife des Messaging-Kanals mit Absender-Autorisierung \\
  \texttt{signal\_\allowbreak dispatcher.service} & Kon\-fi\-gu\-ra\-tion & systemd-Unit der Empfangsschleife \\
  \texttt{notify.py} & Python & Berichts- und Alarmversand mit Stale-Guard \\
  \texttt{notify\_\allowbreak tan.sh} & Shell & photoTAN-Benachrichtigung beim Broker-Login \\
  \texttt{alert\_\allowbreak relay.sh} & Shell & Alarm-Weiterleitung der Datenerfassung an den Messaging-Kanal \\
  \texttt{cleanup\_\allowbreak groups.py} & Python & Werkzeug: Aufräumen verwaister Messaging-Gruppen \\
  \texttt{keepalive.sh} & Shell & Sitzungserhalt der Broker-Session (V2, Overnight-fähig) mit TAN-Fallback \\
  \texttt{keepalive.service} & Kon\-fi\-gu\-ra\-tion & systemd-Unit des Sitzungserhalts \\
  \texttt{keepalive.timer} & Kon\-fi\-gu\-ra\-tion & systemd-Timer des Sitzungserhalts \\
  \texttt{healthcheck.sh} & Shell & Systemüberwachung agent-vm: Dienste, Speicherplatz, Zyklusstatus \\
  \texttt{cycle\_\allowbreak watch.sh} & Shell & Wächter: Alarm bei ausbleibendem Zyklusergebnis \\
  \texttt{heartbeat.sh} & Shell & Heartbeat der Datenerfassung (Lebenszeichen-Datei) \\
  \texttt{heartbeat\_\allowbreak watch.sh} & Shell & Dead-Man-Überwachung des Heartbeats mit Alarmversand \\
  \texttt{loadprobe.sh} & Shell & Lastmessung der agent-vm (Dimensionierungsnachweis) \\
  \texttt{maintain.sh} & Shell & Wartungslauf der Datenerfassung (Rotation, Aufräumen) \\
  \texttt{vacuum.sh} & Shell & SQLite-Vacuum der Marktdatenbank \\
  \texttt{extract\_\allowbreak live\_\allowbreak config.sh} & Shell & Extraktion der Live-Betriebskonfiguration (systemd, Cron, Litestream) von den VMs \\
  \addlinespace[6pt]
  \multicolumn{3}{@{}l}{\textbf{Auswertung}}\\*[2pt]
  \texttt{evaluate\_\allowbreak votes.py} & Python & Wöchentliche Voting-Kalibrierung und ML-Auswertung gegen den Folgetags-Return \\
  \addlinespace[6pt]
  \multicolumn{3}{@{}l}{\textbf{Werkzeuge und Diagnose}}\\*[2pt]
  \texttt{cost\_\allowbreak probe.py} & Python & Ex-Ante-Kostenmessung der Ordervalidierung gegen das Banksystem \\
  \texttt{venue\_\allowbreak probe.py} & Python & Sondierung der Handelsplätze und Orderarten (comdirect) \\
  \texttt{dbg\_\allowbreak positions.py} & Python & Diagnose der Depotbestände \\
  \addlinespace[6pt]
  \multicolumn{3}{@{}l}{\textbf{Dashboard und Journal-Export (webfeed)}}\\*[2pt]
  \texttt{webfeed/\allowbreak export\_\allowbreak journal.py} & Python & Redigierter Journal-Export des Handelsjournals für das Dashboard (collector-vm) \\
  \texttt{webfeed/\allowbreak app/\allowbreak main.py} & Python & Öffentliches Dashboard (FastAPI, Cloud Run) \\
  \texttt{webfeed/\allowbreak app/\allowbreak templates/\allowbreak index.html} & Web & Dashboard-Oberfläche (Diagramme, Journalansicht) \\
  \texttt{webfeed/\allowbreak app/\allowbreak templates/\allowbreak doku.html} & Web & Dokumentationsseite des Dashboards \\
  \texttt{webfeed/\allowbreak app/\allowbreak static/\allowbreak kf-logo.svg} & Web & Logo-Grafik des Dashboards \\
  \texttt{webfeed/\allowbreak app/\allowbreak requirements.txt} & Kon\-fi\-gu\-ra\-tion & Python-Abhängigkeiten des Dashboards \\
  \texttt{webfeed/\allowbreak app/\allowbreak Dockerfile} & Kon\-fi\-gu\-ra\-tion & Container-Bauvorschrift des Dashboards (Cloud Run) \\
  \texttt{webfeed/\allowbreak app/\allowbreak .dockerignore} & Kon\-fi\-gu\-ra\-tion & Build-Ausschlüsse des Container-Images \\
\end{longtable}


\section{Tests}
\label{sec:anhang-archiv-tests}

Die Testsuiten lagen in der Produktivumgebung neben den geprüften Modulen und
sind im Archiv in den Ordner \texttt{Tests/} ausgegliedert. Der Testtreiber
\texttt{run\_tests.sh} stellt den Modulpfad her und führt alle Suiten offline
aus (\texttt{bash Tests/run\_tests.sh}); Suiten mit fehlenden
Drittabhängigkeiten werden dabei als übersprungen ausgewiesen.

%% ---------------------------------------------------------------------------
%% GENERIERT durch build_quellcode_archiv.py am 25.07.2026 08:59 — NICHT von Hand
%% pflegen; Aenderungen in archiv_manifest.csv vornehmen und Skript ausfuehren.
%% ---------------------------------------------------------------------------
\begin{longtable}{@{}p{0.26\textwidth}p{0.26\textwidth}p{0.38\textwidth}@{}}
  \caption[Testsuiten des Archivs]{Ordner \texttt{Tests/} des Archivs: Offline-Testsuiten mit gepr\"uftem Artefakt und Pr\"ufzweck.\quelleeigen}\label{tab:anhang-archiv-tests}\\
  \toprule
  Testdatei & Gepr\"uftes Artefakt & Zweck \\
  \midrule
  \endfirsthead
  \multicolumn{3}{@{}l}{\itshape Fortsetzung von vorheriger Seite}\\
  \toprule
  Testdatei & Gepr\"uftes Artefakt & Zweck \\
  \midrule
  \endhead
  \midrule
  \multicolumn{3}{r@{}}{\itshape Fortsetzung auf der n\"achsten Seite}\\
  \endfoot
  \bottomrule
  \endlastfoot
  \texttt{run\_\allowbreak tests.sh} & alle Testsuiten & Testtreiber: führt alle Suiten mit korrektem Modulpfad (PYTHONPATH) gegen Quellcode/ aus \\
  \texttt{test\_\allowbreak auditor.py} & \texttt{auditor.py} & Multi-LLM-Voting: Mehrheitslogik, Pre-Gate, fail-safe Verhalten (gemockte LLM-Antworten) \\
  \texttt{test\_\allowbreak brokersession.py} & \texttt{broker.py} & Broker-Adapter: Session-, OAuth- und Orderbuchlogik gegen Attrappen \\
  \texttt{test\_\allowbreak control.py} & \texttt{control.py} & Befehlsprotokoll des Steuerkanals einschließlich Absender-Autorisierung \\
  \texttt{test\_\allowbreak credstore.py} & \texttt{credstore.py} & Ver- und Entschlüsselung der Zugangsdatenablage \\
  \texttt{test\_\allowbreak crosscheck.py} & \texttt{datacollect.py} & Zweitquellenprüfung und Sanitätsfilter der Marktdatenerfassung \\
  \texttt{test\_\allowbreak deployments.py} & \texttt{strategyloader.py}, \texttt{agentconfig.yaml} & Konsistenz der Deployment- und Strategiekonfiguration \\
  \texttt{test\_\allowbreak dispatcher.py} & \texttt{signal\_\allowbreak dispatcher.py} & Empfangsschleife des Messaging-Kanals \\
  \texttt{test\_\allowbreak evaluate\_\allowbreak votes.py} & \texttt{evaluate\_\allowbreak votes.py} & Voting- und ML-Auswertung gegen den Folgetags-Return \\
  \texttt{test\_\allowbreak keepalive.sh} & \texttt{keepalive.sh} & Sitzungserhalt: Szenarien einschließlich TAN-Fallback und Nachtruhe \\
  \texttt{test\_\allowbreak mlforecast.py} & \texttt{mlforecast.py} & Statistische Prüfinstanz: Merkmalsbildung, Walk-Forward, Schwellenlogik \\
  \texttt{test\_\allowbreak notify.py} & \texttt{notify.py} & Berichtsversand und Stale-Guard \\
  \texttt{test\_\allowbreak orchestrate.py} & \texttt{orchestrate.py} & Orderausführung: Frische-Gates, Intent-Journal, Reconciliation, Circuit-Breaker \\
  \texttt{webfeed/\allowbreak test\_\allowbreak export.py} & \texttt{webfeed/\allowbreak export\_\allowbreak journal.py} & Journal-Export: Redaktion und Aggregation \\
  \texttt{webfeed/\allowbreak app/\allowbreak test\_\allowbreak app.py} & \texttt{webfeed/\allowbreak app/\allowbreak main.py} & Dashboard-Endpunkte (FastAPI-Testclient) \\
\end{longtable}


\section{Daten}
\label{sec:anhang-archiv-daten}

Die Datenstände belegen die in \cref{sec:umsetzung-evaluation} berichteten
Kennzahlen; ihr jeweiliges Snapshot-Datum ist ausgewiesen. Der Snapshot der
produktiven Marktdatenbank sowie die redigierte Belegkette der
Echtgeld-Orders werden nach Abschluss der Echtgeld-Woche mit der nächsten
Archivversion ergänzt (kursive Einträge).

%% ---------------------------------------------------------------------------
%% GENERIERT durch build_quellcode_archiv.py am 27.07.2026 12:00 — NICHT von Hand
%% pflegen; Aenderungen in archiv_manifest.csv vornehmen und Skript ausfuehren.
%% ---------------------------------------------------------------------------
\begin{longtable}{@{}p{0.28\textwidth}p{0.17\textwidth}p{0.45\textwidth}@{}}
  \caption[Datenst\"ande des Archivs]{Ordner \texttt{Daten/} des Archivs: Daten- und Evaluationsst\"ande mit Snapshot-Datum; kursiv gesetzte Eintr\"age folgen mit der n\"achsten Archivversion nach Abschluss der Echtgeld-Woche.\quelleeigen}\label{tab:anhang-archiv-daten}\\
  \toprule
  Datei & Datenstand & Inhalt \\
  \midrule
  \endfirsthead
  \multicolumn{3}{@{}l}{\itshape Fortsetzung von vorheriger Seite}\\
  \toprule
  Datei & Datenstand & Inhalt \\
  \midrule
  \endhead
  \midrule
  \multicolumn{3}{r@{}}{\itshape Fortsetzung auf der n\"achsten Seite}\\
  \endfoot
  \bottomrule
  \endlastfoot
  \texttt{backtest\_\allowbreak protokoll\_\allowbreak 20260723.txt} & 23.07.2026 & Backtest-Protokoll B0/B1/B2 für 2.000 EUR und 1.000 EUR Kapitaleinsatz, Grundlage der Kennzahlen in Kapitel 3.4 und im Anhang Evaluationsbelege \\
  \texttt{kostentabelle.csv} & 19.07.2026 & Ex-Ante-Transaktionskosten der Bankmessung (Status 201, bankseitig validiert) \\
  \texttt{ml\_\allowbreak walkforward.csv} & 22.07.2026 & Walk-Forward-Kalibrierung der statistischen Prüfinstanz (23.803 bewertete Prognosetage) \\
  \texttt{gspc\_\allowbreak bars1d.csv} & 24.07.2026 & Tagesbars S\&P 500 (Export der Tabelle bars\_1d der Marktdatenbank, ab 30.12.1927) \\
  \texttt{marketdata\_\allowbreak snapshot.sqlite} & \itshape folgt (Live-Woche ab 27.07.2026) & \itshape Snapshot der produktiven Marktdatenbank (collector-vm) zum Abschluss der Echtgeld-Woche \\
  \texttt{journal\_\allowbreak echtgeld\_\allowbreak woche/} & \itshape folgt (Live-Woche ab 27.07.2026) & \itshape Redigierte Journal- und Berichtsauszüge der Echtgeld-Woche (Belegkette der Orders) \\
  \texttt{regime\_\allowbreak profit\_\allowbreak ergebnis.csv} & 24.07.2026 & Ergebnistabelle der Regime-Profitabilitätsanalyse (smatrend/sma über Kapitalstufen, Fenster und Kostenannahmen) — Grundlage der Aussagen zur Handelsgröße \\
  \texttt{veto\_\allowbreak overlay\_\allowbreak ergebnis.csv} & 24.07.2026 & Ergebnistabelle des ML-Veto-Overlays auf den Regime-Arm (Befund negativ) \\
  \texttt{backtest\_\allowbreak v6.csv} & 24.07.2026 & Walk-Forward-Backtest der ML-Prüfinstanz, Konventionsstand v6 \\
  \texttt{mlforecast\_\allowbreak lab\_\allowbreak ergebnis.csv} & 24.07.2026 & Laborergebnisse der Horizont- und Cross-Index-Untersuchung (Zusammenfassung) \\
  \texttt{mlforecast\_\allowbreak lab\_\allowbreak voll.csv} & 24.07.2026 & Laborergebnisse der Horizont- und Cross-Index-Untersuchung (Vollausgabe) \\
  \texttt{lab5\_\allowbreak metriken.csv} & 24.07.2026 & Kennzahlen des Laborlaufs 5 \\
  \texttt{lab5\_\allowbreak tail.csv} & 24.07.2026 & Randverteilung des Laborlaufs 5 \\
  \texttt{lab6\_\allowbreak metriken.csv} & 24.07.2026 & Kennzahlen des Laborlaufs 6 \\
  \texttt{lab6\_\allowbreak tail.csv} & 24.07.2026 & Randverteilung des Laborlaufs 6 \\
  \texttt{orderbuch\_\allowbreak smoketest\_\allowbreak 20260724.csv} & 24.07.2026 & Orderbuchauszug der Orderpfad-Erprobung (Anlage, Änderung, Storno) — Beleg für Statuswerte und Ordermerkmale der Bank \\
  \texttt{venue\_\allowbreak probe2\_\allowbreak protokoll\_\allowbreak 20260724.txt} & 24.07.2026 & Protokoll der Handelsplatz- und Ordertypprüfung inkl. abgelehnter Varianten (HTTP 422) — Beleg für die Grenzen der Bank-API \\
  \texttt{dimensions2\_\allowbreak DE000SB295Z1.json} & 24.07.2026 & Handelsplatz- und Ordertyp-Dimensionen der Bank für das Faktor-Zertifikat — maschineller Beleg zum Protokoll der Handelsplatzprüfung \\
  \texttt{trailing\_\allowbreak validation\_\allowbreak DE000SB295Z1.json} & 24.07.2026 & Validierungsantworten der Bank zu Trailing-/Stop-Varianten inkl. abgelehnter Kombinationen (HTTP 422) \\
\end{longtable}


%  =============================================================================
%  Ende
%  =============================================================================
